With the rapid evolution of LLM agents, long context has emerged as a critical challenge \cite{du-etal-2025-context,bai2024longbench} across open-ended domains like automated software engineering~\citep{xi2023risepotentiallargelanguage}, visual GUI navigation~\citep{nguyen-etal-2025-gui, zhang2025largelanguagemodelbrainedgui}, and deep research~\citep{zhang2025deepresearchsurveyautonomous, huang2025deepresearchagentssystematic}. As agents continuously interact with dynamic environments, their operational history manifests as an \textbf{unbounded agentic trajectory}. This trajectory, typified by the \textbf{interleaved and heterogeneous} ReAct paradigm of Actions, Thoughts, and Observations (A-T-O)~\citep{yao2023reactsynergizingreasoningacting}, rapidly exhausts the LLM's context window to trigger a severe \textbf{\textit{context explosion}}. Such an explosion precipitates a cascade of cognitive failures where information density plummets, critical constraints fade, and long-horizon planning collapses~\citep{liu2023lostmiddlelanguagemodels, Wang_2024,wang-etal-2025-position}. Blindly pursuing longer context windows does not eliminate the need for compression, as models are inherently inept at effectively utilizing vast amounts of uncompressed, noisy textual information~\citep{liu2023lostmiddlelanguagemodels, du-etal-2025-context}, rendering \textbf{effective context management indispensable} for autonomy.

\begin{figure}[t]
  \centering
  \includegraphics[width=\columnwidth]{figures/inrto.pdf}
  \vspace{-.2cm}
  \caption{Illustration of the \textit{context explosion} dilemma and the compression paradigm. Selective distillation transforms unbounded \textcolor{blue}{A-T-O sequences} into \textcolor{red}{structured representations} while preserving operational fidelity.}
  \vspace{-.4cm}
  \label{fig:agent-compression-overview}
\end{figure}



\textbf{Context compression} emerges as a highly effective pathway to mitigate this crisis, exemplified by real-world systems like \raisebox{-0.15em}{
\includegraphics[height=1em]{icons/claude-ai.pdf}}~Claude Code\footnote{\raisebox{-0.15em}{
\includegraphics[height=1em]{icons/claude-ai.pdf}}~Claude Code: https://claude.com/product/claude-code.} that proactively compact contexts upon breaching predefined trajectory thresholds. Unlike the compression of static long documents explored in prior literature~\citep{wang2026fasafrequencyawaresparseattention, li2025surveylargelanguagemodel, xiao2024efficientstreaminglanguagemodels}, the profound \textbf{structural rigidity} of A-T-O sequences in agentic trajectories renders standard text reduction inapplicable. Consequently, as conceptualized in Figure~\ref{fig:agent-compression-overview}, effective agent compression must \textbf{distill massive, noisy trajectories into compact, actionable states} while strictly \textbf{preserving causal action-observation dependencies} and structural integrity.



The prevalence of autonomous agents has catalyzed a rapid proliferation of agentic context compression techniques. These methods are scattered across different compression targets, transformation mechanisms, and control strategies, ranging from pruning and summarization to selective retrieval and representation-level interventions~\citep{lindenbauer2025complexitytrapsimpleobservation, xiao2025improving, wu2026resumunlockinglonghorizonsearch, ye2025agentfoldlonghorizonwebagents, wu2026contextweaverselectivedependencystructuredmemory, wu2026gitcontextcontrollermanage, feng2026agentocrreimaginingagenthistory, wu2026crossmodalmemorycompressionefficient}. Yet these efforts remain largely preliminary and disjointed explorations, lacking any \textbf{systematic synthesis} or \textbf{holistic perspective}. This conceptual fragmentation obscures a principled understanding of both the efficacy of underlying algorithmic designs and how specific compression choices trigger \textbf{distinct downstream agent failures}~\citep{mei2025surveycontextengineeringlarge}. Consequently, practitioners are left without a theoretical framework to navigate the complex trade-offs among \textbf{context footprint}, \textbf{structural fidelity}, and \textbf{historical recoverability}.

Prior surveys on prompt compression~\citep{jiang-etal-2023-llmlingua, li2025prompt} focus on the single-turn reduction of static inputs, whereas we study multi-turn trajectories whose relevance structure evolves over time. Work on KV cache~\citep{li2024survey, liu2025kv, mathurunderstanding} and model compression~\citep{zhu2024survey, liu2025survey, xu2023survey} addresses efficiency at the model or systems layer, while our concern is application-level context control during execution. Surveys on agent memory~\citep{tang2026llm, zhang2025survey, hu2025memory} organize methods by storage form, whereas we emphasize the intervention points, objects, and resulting forms of compressed context. Broader context engineering surveys span retrieval, memory, and orchestration at a high level~\citep{mei2025surveycontextengineeringlarge}; by contrast, we provide a focused treatment of compression as a \textbf{first-class operational capability} in agent systems.

As agents enter the era of long-horizon tasks, context compression has become a three-dimensional systems problem rather than a single-turn prompt optimization trick. Yet a systematic review of this emerging domain remains absent. To fill this gap, we introduce a unified taxonomy organized around three questions: \textbf{\textit{compression target}} (what is compressed), \textbf{\textit{compression mechanism}} (how it is transformed and retained), and \textbf{\textit{control policy}} (who decides when compression is triggered). These dimensions map directly onto the operational pipeline: targets specify \textbf{Select}, mechanisms instantiate \textbf{Compress} and \textbf{Store}, and policies govern when these stages, and when needed \textbf{Recover}, are activated. We further formalize three recurring failure modes, \fone, \ftwo, and \fthree, and use them to expose reliability trade-offs. Finally, we critique current evaluation paradigms~\citep{liu2025agentbenchevaluatingllmsagents} and argue for metrics beyond end-task success, including information density, temporal consistency, and error propagation.

The rest of the survey is organized as follows. Section~\ref{sec:2_background} introduces the background and unified pipeline, Section~\ref{sec:3_formal_view} presents a formal lens, Sections~\ref{sec:3_taxonomy} and~\ref{sec:4_failure_modes} develop the taxonomy and failure analysis, and the remaining sections discuss domain-specific trade-offs, evaluation, future directions, and conclusion.
